\documentclass[11pt,letterpaper]{article}
\usepackage[T1]{fontenc}
\usepackage[utf8]{inputenc}
\usepackage{amsmath,amsthm}
\usepackage{newpxtext,newpxmath}
\usepackage[margin=1in,headheight=15pt]{geometry}
\usepackage{microtype,mathtools}
\usepackage{booktabs,array,longtable,enumitem}
\usepackage[dvipsnames]{xcolor}
\usepackage[most]{tcolorbox}
\usepackage{fancyhdr,titlesec,needspace,aliascnt}
\usepackage{listings}
\usepackage{hyperref}
\usepackage[nameinlink,noabbrev]{cleveref}
\definecolor{Ink}{HTML}{18354A}
\definecolor{Accent}{HTML}{176E74}
\definecolor{Pale}{HTML}{EFF6F6}
\definecolor{Muted}{HTML}{52616B}
\hypersetup{colorlinks=true,linkcolor=Accent,citecolor=Accent,urlcolor=Accent,
 pdftitle={Friendly Order Type and Incomparability Components},
 pdfauthor={ChatGPT},pdfsubject={Finite posets, ordinal invariants, multiset ordering}}
\titleformat{\section}{\Large\bfseries\color{Ink}}{\thesection}{0.65em}{}
\titleformat{\subsection}{\large\bfseries\color{Ink}}{\thesubsection}{0.6em}{}
\pagestyle{fancy}\fancyhf{}
\fancyhead[L]{\small\color{Muted}Friendly order type and incomparability components}
\fancyhead[R]{\small\color{Muted}Research report}
\fancyfoot[C]{\small\thepage}
\renewcommand{\headrulewidth}{0.3pt}
\setlength{\parskip}{4pt}
\setlength{\parindent}{1.25em}
\setlist{itemsep=2pt,topsep=5pt}
\newtheorem{theorem}{Theorem}[section]
\newaliascnt{lemma}{theorem}
\newtheorem{lemma}[lemma]{Lemma}
\aliascntresetthe{lemma}
\crefname{lemma}{lemma}{lemmas}
\Crefname{lemma}{Lemma}{Lemmas}
\newaliascnt{proposition}{theorem}
\newtheorem{proposition}[proposition]{Proposition}
\aliascntresetthe{proposition}
\crefname{proposition}{proposition}{propositions}
\Crefname{proposition}{Proposition}{Propositions}
\newaliascnt{corollary}{theorem}
\newtheorem{corollary}[corollary]{Corollary}
\aliascntresetthe{corollary}
\crefname{corollary}{corollary}{corollarys}
\Crefname{corollary}{Corollary}{Corollarys}
\theoremstyle{definition}
\newaliascnt{definition}{theorem}
\newtheorem{definition}[definition]{Definition}
\aliascntresetthe{definition}
\crefname{definition}{definition}{definitions}
\Crefname{definition}{Definition}{Definitions}
\newaliascnt{example}{theorem}
\newtheorem{example}[example]{Example}
\aliascntresetthe{example}
\crefname{example}{example}{examples}
\Crefname{example}{Example}{Examples}
\newaliascnt{problem}{theorem}
\newtheorem{problem}[problem]{Problem}
\aliascntresetthe{problem}
\crefname{problem}{problem}{problems}
\Crefname{problem}{Problem}{Problems}
\theoremstyle{remark}
\newaliascnt{remark}{theorem}
\newtheorem{remark}[remark]{Remark}
\aliascntresetthe{remark}
\crefname{remark}{remark}{remarks}
\Crefname{remark}{Remark}{Remarks}
\newcommand{\fot}{\operatorname{f}}
\newcommand{\Inc}{\operatorname{Inc}}
\newcommand{\cc}{\operatorname{c}}
\newcommand{\Bad}{\operatorname{Bad}}
\newcommand{\Ant}{\operatorname{Ant}}
\newcommand{\rk}{\operatorname{rk}}
\newcommand{\Min}{\operatorname{Min}}
\newcommand{\str}{\operatorname{str}}
\newcommand{\DM}{\mathrm{DM}}
\newcommand{\M}{\mathcal M}
\newcommand{\N}{\mathbb N}
\newcommand{\C}{\mathsf C}
\newcommand{\A}{\mathsf A}
\newcommand{\epsbot}{\varepsilon_{\bot}}
\newcommand{\epstop}{\varepsilon_{\top}}
\newtcolorbox{resultbox}{colback=Pale,colframe=Accent,boxrule=0.7pt,
 arc=2pt,left=10pt,right=10pt,top=8pt,bottom=8pt,breakable}
\lstset{basicstyle=\small\ttfamily,breaklines=true,columns=fullflexible,
 frame=single,rulecolor=\color{Muted},backgroundcolor=\color{Pale},
 showstringspaces=false,aboveskip=9pt,belowskip=9pt}
\begin{document}
\begin{titlepage}
\vspace*{0.5in}
{\small\sffamily\bfseries\color{Accent}ORDER THEORY \; / \; ORDINAL INVARIANTS}\par
\vspace{0.3in}
{\Huge\bfseries\color{Ink}Friendly Order Type and\par
Incomparability Components}\par
\vspace{0.22in}
{\Large A finite-case answer to a compositional question,\par
with transfinite extensions and executable certificates}\par
\vspace{0.35in}
{\large 19 September 2026}\par
\vspace{0.15in}
{\color{Muted}An AI-assisted mathematical research report prepared for\par
Vladimir Reshetnikov}\par
\vspace{0.35in}
\begin{resultbox}
For every finite poset $P$,
\[
\boxed{\fot(P)=|P|-\cc\bigl(\Inc(P)\bigr).}
\]
Here $\fot$ is Vialard's friendly order type and $\cc$ counts all connected
components of the incomparability graph, including isolated vertices.
Consequently,
\[
\boxed{w\bigl(\M_{\DM}(P)\bigr)=\omega^{\,|P|-\cc(\Inc(P))}.}
\]
The report proves the component formula, classifies maximum friendly subsets,
and gives composition laws, transfinite examples, and independently checked
finite computations.
\end{resultbox}
\vfill
\noindent\textbf{Scope and status.}
The source problem asks for a broader understanding and compositional
calculus of friendly order type. This report resolves the finite specialization
and a class of well-ordered sums; it does not resolve that entire program.
The arguments are presented as proofs, but the manuscript has not been
independently refereed or formally verified. A literature search did not locate
the component formula; that is not a certification of priority.
\end{titlepage}

\begin{abstract}
Friendly order type is the rank of the tree of bad sequences whose successive
choices retain an incomparable witness in the current residual. It was
introduced to compute the ordinal width of the multiset ordering. We prove
that for a finite poset it is simply the rank of the graphic matroid of the
incomparability graph: the number of vertices minus the number of components.
The key lemma produces a maximal non-cut vertex while preserving any specified
minimal vertex. It yields optimal sequences, spanning-forest certificates, and
a classification of all maximum friendly subsets. Further consequences include
closed formulas for finite Cartesian products and arbitrary nonempty finite
lexicographic substitutions, an exact generating function for the distribution
on naturally labelled posets, and a transfinite summation theorem. In particular,
well-ordered sums of two-element antichains realize every ordinal as a friendly
order type, even though all countably infinite examples have isomorphic
incomparability graphs. Computations compare the defining recursion with the
formula on 101,660 naturally labelled posets through size seven, and test
additional product, substitution, and prescribed-root cases.
\end{abstract}
\setcounter{tocdepth}{1}
\tableofcontents
\newpage

\section{The problem selected and the contribution}
\label{sec:problem}
The question selected for this investigation is the compositional study of
\emph{friendly order type}. Vialard's thesis explicitly asks whether this
invariant can be calculated compositionally and whether it has a further
structural interpretation \cite[Conclusion, p.~109]{VialardThesis}.
The finite specialization is particularly concrete: its definition is a
recursive search over residual posets, but one can ask for a direct structural
formula that avoids this search altogether.

\begin{problem}[Finite specialization investigated here]
Given an arbitrary finite poset $P$, identify a nonrecursive structural
expression for its friendly order type. Use that expression to obtain
composition laws and explicit optimal witnesses.
\end{problem}

This is our specialization of the published research direction, not a claim
that the literature separately labels this exact finite statement as a
conjecture. There is already a substantial theory: the invariant and its
connection to multiset width are due to Vialard, and the ordinal-sum law is
also known \cite[Definition~3.3, Theorem~3.4, Proposition~4.1(1)]{VialardMFCS}.
Those facts must not be represented as discoveries of this report.

The proposed contribution is a complete finite structural description,
together with consequences proved below. The proof uses two ideas. First,
incomparability components are not arbitrary pieces of a poset: they sit in
a linearly ordered sequence. Second, in a finite connected incomparability
graph one can remove a \emph{maximal poset element} without disconnecting
the graph. Repeating this operation leaves exactly one element per component.
The defining sequence process can select every other element and no more.

The finite invariant therefore has a particularly simple graph-theoretic
interpretation. Its transfinite analogue does not forget the ordering of the
components; this distinction will be made explicit rather than hidden in an
informal replacement of finite numbers by ordinals.

\subsection{Literature and novelty audit}
The audit used the published MFCS 2023 article, the author's 2024 thesis,
the earlier arXiv version, and the author's publication page, accessed on
19 September 2026 \cite{VialardMFCS,VialardThesis,VialardArxiv,VialardHome}.
The earlier preprint uses the terminology \emph{maximal safe order type};
this report uses the later tree-based definition of \emph{friendly order type}.
Exact-phrase searches for that terminology and for an incomparability-component
formula did not locate a later treatment of the formula proved here.

This audit establishes the provenance of the research question, but does not
prove that the formula is unpublished. Its elementary ingredients may be
known in other terminology. The mathematical assertions below have explicit
proofs independent of any assertion of novelty. General infinite formulas
not used in those proofs are deliberately not imported from earlier source
versions.

\section{Definitions, residuals, and rank conventions}
\label{sec:definitions}
We work in ZFC with set-sized orders. A \emph{poset} is a set with a reflexive, antisymmetric, transitive relation
$\leq$. We write $x\parallel y$ when neither $x\leq y$ nor $y\leq x$.
The \emph{incomparability graph} $\Inc(P)$ has vertex set $P$ and an
undirected edge between distinct incomparable elements. Its number of
connected components is denoted by $\cc(\Inc(P))$; an isolated vertex counts
as one component and the empty graph has zero components.

A finite sequence $s=(x_0,\ldots,x_{k-1})$ is \emph{bad} if
$x_i\not\leq x_j$ whenever $i<j$. The residual after this sequence is
\[
 P_s=P\setminus\bigcup_{i<k}\mathord\uparrow x_i,
 \qquad \mathord\uparrow x=\{y\in P:x\leq y\}.
\]
Thus the next member of a bad sequence must lie in $P_s$. A poset is a
\emph{well-partial order}, or wpo, if it has no infinite bad sequence.
Every finite poset is a wpo.

The tree of finite bad sequences, rooted at the empty sequence, is
well-founded for a wpo. For a node $s$ in any such tree, our rank convention is
\begin{equation}\label{eq:tree-rank}
 \rk(s)=\sup\{\rk(sx)+1:sx\text{ is an immediate extension of }s\},
 \qquad \sup\varnothing=0.
\end{equation}
The rank of a tree means the rank of its root. In a finite tree this is the
maximum number of edges on a root-to-leaf path. These are the standard
rank conventions for ordinal invariants of wpos \cite{DSS}.

\begin{definition}[Friendly order type]
An \emph{open-ended bad sequence} is obtained by repeatedly choosing
$x\in P_s$ for which some $y\in P_s$ satisfies $x\parallel y$.
The element $y$ is a \emph{friend} of that choice. The empty sequence is
allowed. The friendly order type $\fot(P)$ is the rank of the tree of these
sequences. This is the invariant denoted $o_\perp(P)$ by Vialard
\cite[Definition~3.3]{VialardMFCS}.
\end{definition}

The friend is required in the residual \emph{before} the choice is made.
Since it is incomparable to the choice, it also survives that individual
choice. It need not survive all later choices, and it need not eventually
be selected. These distinctions matter: requiring a permanent friend or
requiring a friend in the original poset would define different processes.

Write $P\perp x=\{y\in P:y\parallel x\}$ and
$P\not\geq x=P\setminus\mathord\uparrow x$. Directly from the definition,
\begin{equation}\label{eq:residual}
 \fot(P)=\sup_{\substack{x\in P\\P\perp x\ne\varnothing}}
 \bigl(\fot(P\not\geq x)+1\bigr).
\end{equation}
For finite $P$ the supremum is a maximum. The next small observations also
apply to infinite wpos.

\begin{lemma}\label{lem:elementary}
If $Q$ is an induced subposet of a wpo $P$, then $\fot(Q)\leq\fot(P)$.
Moreover, $\fot(P)=0$ exactly when $P$ is a chain, including the empty chain.
Every subsequence of an open-ended bad sequence is open-ended.
\end{lemma}
\begin{proof}
A sequence using only elements of $Q$ has the same comparisons in $P$.
Each friend available inside $Q$ is still available inside $P$, so its tree
embeds as a subtree and ranks cannot increase on restriction. A first move
exists precisely when some incomparable pair exists. Finally, deleting
previous choices from a sequence only enlarges each subsequent residual.
The friends witnessing the original sequence therefore witness any
subsequence.
\end{proof}

For a finite nonempty poset with $n$ elements one immediately has
\begin{equation}\label{eq:rough-bound}
 0\leq\fot(P)\leq n-1.
\end{equation}
Indeed, all selected elements are distinct, and the last selected element
needs a different element still present as a friend. Our task is to determine
exactly how much stronger this bound becomes when the incomparability graph
is disconnected.

\section{The order structure of incomparability components}
\label{sec:components}
For posets $A,B$, their \emph{ordinal sum} $A+B$ retains the internal orders
and places every element of $A$ below every element of $B$. An ordinal sum
is very different from a disjoint union, which adds no cross comparisons.

\begin{lemma}[Ordered components]\label{lem:ordered-components}
Distinct connected components of $\Inc(P)$ are uniformly comparable:
for two such components $C,D$, either every member of $C$ is below every
member of $D$, or conversely. These directions linearly order the components.
If $P$ is a wpo, this component order is a well-order.
\end{lemma}
\begin{proof}
Elements in different components cannot be incomparable, since that would
supply an edge joining the components. Suppose $a\parallel b$ and $z$ lies
outside their component. If $a<z$, then $z<b$ would imply $a<b$, a
contradiction. Since $b,z$ are comparable, necessarily $b<z$. The same
argument handles the opposite direction.

Propagating along an incomparability path shows that the direction between
$z$ and all vertices of a component is constant. Propagating within the
other component makes the direction uniform on both sides. Transitivity of
the poset then gives transitivity of the component order, and distinct
components are comparable, so it is linear.

If the component order of a wpo had an infinite strictly descending sequence,
choosing one element from each component would give an infinite strictly
descending sequence in $P$. This is impossible. In the usual set-theoretic
setting, a linear order with no infinite descending sequence is a well-order.
\end{proof}

In particular, every finite poset has a canonical decomposition
\begin{equation}\label{eq:component-sum}
 P=C_1+\cdots+C_c,
\end{equation}
where the $C_i$ are its incomparability components in increasing order.
This is a decomposition into \emph{induced posets}; it does not assert that
the components are chains or antichains.

\begin{lemma}[Component capacity]\label{lem:capacity}
In any open-ended bad sequence in a finite poset, at most $|C|-1$ elements
can be selected from each incomparability component $C$.
\end{lemma}
\begin{proof}
A friend of an element of $C$ must also lie in $C$. If a sequence selected
every element of $C$, consider its last choice from $C$. All other elements
of $C$ would already have been selected and hence removed. No friend could
remain, a contradiction. Extra removals caused by upward closures cannot
repair this obstruction.
\end{proof}

\begin{corollary}\label{cor:upper}
For every finite poset,
\[
 \fot(P)\leq |P|-\cc(\Inc(P)).
\]
\end{corollary}

The main issue is the lower bound: can these local capacities all be reached
simultaneously? A generic graph deletion order does not suffice. To avoid
accidentally removing further elements, the next selected vertex should be
maximal in the remaining poset. The following lemma supplies exactly this.

\section{A maximal non-cut lemma and the finite formula}
\label{sec:main}
A vertex is called \emph{non-cut} here if deleting it leaves a connected graph.
We apply this only to connected graphs with at least two vertices; the
one-vertex graph left after a deletion is connected.

\begin{lemma}[Maximal non-cut vertex with a protected root]\label{lem:noncut}
Let $P$ be a finite poset with at least two elements and connected
incomparability graph. Given any minimal element $r$ of $P$, there exists
$x\ne r$ such that $x$ is maximal in $P$ and $\Inc(P\setminus\{x\})$
is connected.
\end{lemma}
\begin{proof}
Choose a maximal element $x\ne r$. Such an element exists: if the only
maximal element were $r$, finiteness would imply that every element is
below $r$, contrary to the minimality of $r$ and $|P|\geq2$.
If deleting $x$ leaves a connected graph, we are done.

Otherwise let the components of $\Inc(P\setminus\{x\})$ be
\[
 D_1<D_2<\cdots<D_k,\qquad k\geq2,
\]
in the order provided by \cref{lem:ordered-components}. Each $D_i$ has a
vertex adjacent to $x$ in the original incomparability graph. Otherwise
$D_i$ would have no edge either to $x$ or to any other $D_j$, contradicting
connectedness of the original graph.

Choose $z\in D_1$ with $z\parallel x$. For any $i\geq2$ and $y\in D_i$,
we have $z<y$. If $y<x$, transitivity would give $z<x$, a contradiction.
If $x<y$, maximality of $x$ would be contradicted. Hence
\begin{equation}\label{eq:universal-x}
 x\parallel y\qquad\text{for every }y\in D_2\cup\cdots\cup D_k.
\end{equation}
Choose a maximal element $y$ of the last component $D_k$. It is maximal
in all of $P$: any element in an earlier component is below it, no element
of $D_k$ is above it, and it is incomparable with $x$.
Also $y\ne r$, because any element of $D_1$ is below $y$, while $r$ is
minimal.

Deleting $y$ leaves a connected graph. Indeed, $D_1$ remains connected and
has an edge to $x$. Every remaining vertex of $D_i$ for $i\geq2$ is directly
adjacent to $x$ by \eqref{eq:universal-x}, regardless of whether deletion of
$y$ disconnects the graph inside $D_k$. Thus $y$ is the required vertex
when the original choice $x$ is a cut vertex.
\end{proof}

\begin{proposition}[Connected case]\label{prop:connected}
If $P$ is finite and $\Inc(P)$ is connected, then
\[
 \fot(P)=|P|-1.
\]
More strongly, for every minimal $r\in P$ there is an open-ended sequence
selecting precisely $P\setminus\{r\}$. At each step it selects a maximal
vertex of the current induced poset and leaves that poset connected.
\end{proposition}
\begin{proof}
A singleton gives the empty sequence. For at least two elements, apply
\cref{lem:noncut} with protected root $r$. The chosen $x$ has a friend since
the graph is connected and has more than one vertex. Since $x$ is maximal,
its upward closure in the current poset is just $\{x\}$; hence the residual
is exactly $P\setminus\{x\}$, with a connected incomparability graph.
The root $r$ is still minimal there. Induction repeats this construction
until only $r$ remains. The sequence has $|P|-1$ choices, and
\eqref{eq:rough-bound} proves optimality.
\end{proof}

\Needspace{9\baselineskip}
\begin{theorem}[Finite component formula]\label{thm:main}
For every finite poset $P$, including the empty poset,
\begin{equation}\label{eq:main}
 \boxed{\fot(P)=|P|-\cc(\Inc(P)).}
\end{equation}
\end{theorem}
\begin{proof}
The empty case is immediate. In the decomposition
$P=C_1+\cdots+C_c$, choose a minimal root $r_i$ in each induced component
$C_i$. By \cref{prop:connected}, there is a local open-ended sequence of
length $|C_i|-1$ leaving only $r_i$ in that component.

Run these local sequences in the order $C_c,C_{c-1},\ldots,C_1$.
A choice in one component cannot remove an element of an earlier component,
so a component is still intact when its local sequence starts. A choice may
remove surviving roots in later components, but those roots are not needed
as friends: all local friends lie in the component currently being processed.
Thus the concatenated sequence is legal and has length
\[
 \sum_{i=1}^c(|C_i|-1)=|P|-c.
\]
The upper bound is \cref{cor:upper}.
\end{proof}

\begin{remark}[Why the proof is not just a graph argument]
For arbitrary graphs one can repeatedly delete non-cut vertices. The extra
requirement here is maximality in a compatible poset. Without it, deleting
an upward closure could remove several vertices at once and destroy the
intended sequence length. \Cref{lem:noncut} is the point where transitivity
of the order supplies more structure than graph connectivity alone.
\end{remark}

\section{All maximum friendly subsets and forest certificates}
\label{sec:subsets}
For a finite subset $S\subseteq P$, say that $S$ is \emph{friendly in $P$}
if it admits a linear extension whose reverse listing is an open-ended bad
sequence in $P$. By \cref{lem:elementary}, all subsequences of that listing
are then open-ended as well. Conversely, an open-ended sequence listing
all elements of $S$ has a reverse listing that is a linear extension of $S$.
This finite formulation agrees with the friendly-subset condition used in
Vialard's alternative characterization \cite[Definition~4.4]{VialardMFCS}.

\begin{theorem}[Classification of maximum friendly subsets]\label{thm:subsets}
Let $C_1<\cdots<C_c$ be the incomparability components of a finite poset $P$.
The friendly subsets of maximum possible size are exactly
\begin{equation}\label{eq:subset-class}
 P\setminus\{r_1,\ldots,r_c\},\qquad r_i\in\Min(C_i).
\end{equation}
Their common size is $|P|-c$, and their number is
\begin{equation}\label{eq:subset-count}
 \prod_{i=1}^c |\Min(C_i)|,
\end{equation}
with empty product equal to $1$.
\end{theorem}
\begin{proof}
The construction in \cref{thm:main} realizes every choice of roots in
\eqref{eq:subset-class}, proving sufficiency and the asserted size.
For necessity, a maximum friendly subset must omit at least one member
of each component by \cref{lem:capacity}. Since its size is $|P|-c$, it
omits exactly one, say $r_i$, from each $C_i$.

Fix a sequence listing the subset, and consider a component $C_i$.
If $r_i$ were not minimal in $C_i$, some selected element $a\in C_i$
would satisfy $a<r_i$. At the choice of $a$, the root $r_i$ would be removed.
If this were not the last choice from $C_i$, then no unselected member of
$C_i$ could remain as a friend for the last choice. If it were the last
choice, the only possibly remaining other member of $C_i$ would be $r_i$,
which is comparable to $a$. Both possibilities violate the friend condition.
Thus each $r_i$ is minimal. Different root choices give different complements,
so multiplying the numbers of choices yields \eqref{eq:subset-count}.
\end{proof}

The theorem is stronger than merely computing a rank. It identifies every
largest witness subset, although it does not enumerate all permissible
orderings of a given witness subset.

\begin{proposition}[Spanning-forest certificate]\label{prop:forest}
The construction in \cref{thm:main} can attach to every choice $x$ a friend
$y$ such that the edges $\{x,y\}$ form a spanning tree in each original
incomparability component. The unselected $r_i$ are the roots of those trees.
\end{proposition}
\begin{proof}
Within a component, choose as friend any neighbor remaining after deleting
the selected maximal vertex. Orient the witness edge from that vertex to
its friend. The endpoint is either selected later in that local sequence
or is its protected root. Following oriented edges therefore strictly
increases deletion time until reaching the root; cycles are impossible.
Every non-root has exactly one outgoing edge and eventually reaches the
root, so the edges form a connected acyclic spanning graph. The construction
never uses a friend from a different component, giving a spanning forest
on the whole poset.
\end{proof}

\begin{corollary}[Graphic matroid rank and finite duality]\label{cor:graphic}
For finite $P$, $\fot(P)$ is the maximum number of edges in a forest of
$\Inc(P)$. In particular,
\[
 \fot(P^{\mathrm{op}})=\fot(P).
\]
\end{corollary}
\begin{proof}
A tree on $m$ vertices has $m-1$ edges, so a maximum forest in a finite graph
has one spanning tree per component and $|P|-c$ edges. This is precisely
the rank of its graphic matroid. Reversing the order does not change which
pairs are incomparable, so it does not change this rank.
\end{proof}

Not every spanning forest, and not every graph deletion order, is automatically
a legal certificate. The theorem concerns the value of the graphic rank;
\cref{prop:forest} supplies particular forests compatible with the residual
process.

\section{Consequences for the ordinal width of multiset orders}
\label{sec:multisets}
A finite multiset on $P$ is a finitely supported function $P\to\N$.
After cancelling the common part of two multisets $\mu,\nu$, the strict
multiset ordering $\mu<_{\DM}\nu$ holds if the remainder of $\nu$ is
nonempty and every element remaining in $\mu$ is strictly below some
remaining element of $\nu$. A witness in $\nu$ may be reused. We add equality
to obtain $\leq_{\DM}$. This is the multiset ordering, not the injective
multiset embedding order \cite[Section~1]{VialardMFCS}.

For a wpo $Q$, its \emph{ordinal width} $w(Q)$ is the rank of the tree of
finite sequences of pairwise incomparable elements, with convention
\eqref{eq:tree-rank}. This agrees with ordinary antichain width on finite
posets, but may be a transfinite ordinal on an infinite wpo.

\begin{theorem}[Vialard's width transfer, imported]\label{thm:transfer}
For every wpo $P$,
\begin{equation}\label{eq:transfer}
 w\bigl(\M_{\DM}(P)\bigr)=\omega^{\fot(P)}.
\end{equation}
\end{theorem}
This is the existing multiset-width theorem
\cite[Theorem~5.3.3]{VialardThesis}; its technical proof is not reproduced
or claimed as new here. All the finite computations of its exponent below
follow from the preceding self-contained argument.

\begin{corollary}[Closed multiset-width formula]\label{cor:multiset}
For every finite poset $P$,
\[
 w\bigl(\M_{\DM}(P)\bigr)=\omega^{\,|P|-\cc(\Inc(P))}.
\]
\end{corollary}
\begin{proof}
Substitute \cref{thm:main} into \eqref{eq:transfer}.
\end{proof}

The case $P=\varnothing$ is worth keeping: there is exactly one multiset,
the empty one, so the width is $1=\omega^0$, not $0$.
For a nonempty finite chain the exponent is also zero, consistently with
its multiset ordering being a chain.

\begin{example}[The three-element distinction]\label{ex:three}
Let $V=\C_1+\A_2$ be a bottom element below two incomparable maxima, and let
$D=\C_2\sqcup\C_1$ be a two-element chain disjoint from a singleton.
Here $\C_m$ denotes an $m$-element chain and $\A_m$ an $m$-element antichain.
Both posets have cardinality and maximal order type $3$, height $2$, and
ordinary width $2$. But $\Inc(V)$ has two components whereas $\Inc(D)$ is
connected. Therefore
\[
 \fot(V)=1,\quad \fot(D)=2,\qquad
 w(\M_{\DM}(V))=\omega,\quad w(\M_{\DM}(D))=\omega^2.
\]
Thus the failure of determination by cardinality, height, and ordinary width
already occurs on three elements. On at most two elements, those data
distinguish the possible isomorphism types, so this size is minimal.
\end{example}

\section{Finite composition laws}
\label{sec:composition}
All the formulas in this section apply to arbitrary finite posets, not only
to posets generated from chains by a restricted list of operations.

\subsection{Ordinal sums and disjoint unions}
\begin{corollary}\label{cor:finite-sums}
For finite $A,B$,
\[
 \fot(A+B)=\fot(A)+\fot(B).
\]
If $A,B$ are both nonempty, then
\[
 \fot(A\sqcup B)=|A|+|B|-1.
\]
\end{corollary}
\begin{proof}
The graph of an ordinal sum is the disjoint union of the two incomparability
graphs, so both vertex and component counts add. In the poset disjoint union,
every vertex of $A$ is incomparable to every vertex of $B$. The resulting
graph is connected when both are nonempty. Apply \cref{thm:main}.
\end{proof}
The first identity is a finite instance of a known ordinal-sum law, and the
finite disjoint-union identity is already consistent with the published
finite formulas. The contribution here is their common structural explanation,
not priority for these special cases.

\subsection{Arbitrary lexicographic substitution}
Let $P$ be a finite base poset. For every $p\in P$, let $Q_p$ be a nonempty
finite poset. The \emph{lexicographic substitution} $P[Q_p]$ consists of
pairs $(p,q)$ with $q\in Q_p$, ordered by
\[
 (p,q)\leq(p',q')
 \quad\Longleftrightarrow\quad
 p<p'\ \text{ or }\ (p=p'\text{ and }q\leq q').
\]
In particular, fibers over incomparable base elements are completely
incomparable to one another.

\begin{theorem}[Substitution formula]\label{thm:substitution}
Let $\mathcal C$ be the set of components of $\Inc(P)$. Then
\begin{equation}\label{eq:substitution}
 \boxed{\displaystyle
 \fot(P[Q_p])=
 \sum_{\substack{C\in\mathcal C\\|C|\geq2}}
       \left(\sum_{p\in C}|Q_p|-1\right)
 +\sum_{\{p\}\in\mathcal C}\fot(Q_p).}
\end{equation}
\end{theorem}
\begin{proof}
There are no incomparability edges between fibers lying over different base
components. Over a singleton component $\{p\}$, the incomparability graph
is just $\Inc(Q_p)$ and its contribution is $\fot(Q_p)$.

Now let $C$ be a connected base component with at least two vertices. Each
base vertex has a neighbor, and fibers are nonempty. Any two vertices in
the same fiber can be joined by a path of length two through a neighboring
fiber. A path in $\Inc(P)$ connecting two base vertices lifts to a path
connecting any prescribed vertices in their fibers, since adjacent fibers
are joined by all possible edges. Thus the entire collection of fibers over
$C$ is one connected component, with $\sum_{p\in C}|Q_p|$ vertices.
Its friendly contribution is that number minus one. Summing proves the formula.
\end{proof}

\begin{corollary}[Uniform fibers]\label{cor:uniform}
Suppose all fibers are the same nonempty finite poset $Q$ of size $m$.
If $P$ has $n$ elements, $s$ singleton incomparability components, and $t$
non-singleton incomparability components, then
\[
 \fot(P[Q])=m(n-s)-t+s\fot(Q).
\]
\end{corollary}

Empty fibers cannot simply be substituted into \eqref{eq:substitution}:
they may destroy paths in the base incomparability graph. When some fibers
are empty, first delete their base vertices, take the induced base order,
and apply the theorem to the remaining nonempty fibers.

\subsection{Cartesian products}
For a nonempty finite poset $P$, put
\[
 \epsbot(P)=
 \begin{cases}1&\text{if $P$ has a least element},\\0&\text{otherwise},\end{cases}
 \quad
 \epstop(P)=
 \begin{cases}1&\text{if $P$ has a greatest element},\\0&\text{otherwise}.\end{cases}
\]
These are global least and greatest elements, not merely minimal and
maximal elements.

\begin{theorem}[Cartesian-product formula]\label{thm:cartesian}
Let $A,B$ be finite posets with $|A|,|B|\geq2$. Then
\begin{equation}\label{eq:cartesian}
 \boxed{\fot(A\times B)=|A||B|-1
 -\epsbot(A)\epsbot(B)-\epstop(A)\epstop(B).}
\end{equation}
Equivalently, the incomparability graph has exactly one non-singleton
component, together with an isolated joint least element when both factors
have least elements and an isolated joint greatest element when both have
greatest elements.
\end{theorem}
\begin{proof}
Choose linear extensions of $A$ and $B$, listing them as
$a_0,\ldots,a_{m-1}$ and $b_0,\ldots,b_{n-1}$. Adding comparisons to make
each list a chain can only delete incomparability edges from the product.
Consequently the incomparability graph of the chain grid
$\C_m\times\C_n$ is a spanning subgraph of $\Inc(A\times B)$.

In that grid, all vertices except $(0,0)$ and $(m-1,n-1)$ form a connected
incomparability graph. To see this, use the two incomparable anchors
\[
 u=(0,n-1),\qquad v=(m-1,0).
\]
A vertex $(i,j)$ with $i>0$ and $j<n-1$ is adjacent to $u$; a vertex with
$i<m-1$ and $j>0$ is adjacent to $v$. Every non-corner vertex is an anchor
or satisfies at least one of these conditions. The anchors are adjacent,
so these vertices form one connected component in the original product too.

It remains to classify the two possible corner exceptions. The lower corner
$(a_0,b_0)$ is isolated if $a_0,b_0$ are both least elements, since it is
then below every vertex. Otherwise, say $a_0$ is not least. Being first in
a linear extension, it is minimal, so some $a\in A$ is incomparable with
$a_0$. Then $(a,b_0)$ is incomparable with the lower corner. It is a
non-corner grid vertex because $a\ne a_0$ and $b_0\ne b_{n-1}$.
Thus the lower corner joins the large component unless it is genuinely a
least element. The argument for the upper corner is dual.

No other component can exist. The number of components is therefore
$1+\epsbot(A)\epsbot(B)+\epstop(A)\epstop(B)$, and \cref{thm:main} gives
the formula.
\end{proof}

If a factor is empty, the product is empty. If one factor is a singleton,
the product is isomorphic to the other factor; these cases must not be put
into \eqref{eq:cartesian} without its size hypotheses.

\begin{corollary}[Several factors and Boolean lattices]\label{cor:boolean}
For finite $P_1,\ldots,P_d$ with $d\geq2$ and $|P_i|\geq2$,
\[
 \fot\left(\prod_{i=1}^dP_i\right)
 =\prod_{i=1}^d|P_i|-1
  -\prod_{i=1}^d\epsbot(P_i)-\prod_{i=1}^d\epstop(P_i).
\]
In particular, for the Boolean lattice $B_d$ with $d\geq2$,
\[
 \fot(B_d)=2^d-3,
 \qquad w\bigl(\M_{\DM}(B_d)\bigr)=\omega^{2^d-3}.
\]
\end{corollary}
\begin{proof}
Apply \cref{thm:cartesian} successively; a product has a least, respectively
greatest, element exactly when every factor does. Now use
$B_d\cong\C_2^d$ and \cref{cor:multiset}.
\end{proof}

\section{Well-ordered sums and genuine transfinite values}
\label{sec:transfinite}
The finite graph formula involves ordinary subtraction. There is no reason
to replace it by an unspecified ordinal subtraction. The correct extension
on well-ordered sums uses \emph{ordinal addition}.

\begin{lemma}[Binary ordinal sums]\label{lem:binary-ordinal}
For arbitrary wpos $A,B$,
\begin{equation}\label{eq:binary-ordinal}
 \fot(A+B)=\fot(A)+\fot(B),
\end{equation}
where the addition on the right is ordinal addition.
\end{lemma}
\begin{proof}
We give a residual proof of the known sum law. Fix $A$ and argue by induction
on the rank of $\Bad(B)$. A move from $A$ has friends only in $A$ and leaves
$A\not\geq x$, deleting all of $B$. A move from $B$ has friends only in
$B$ and leaves $A+(B\not\geq x)$. The latter residual of $B$ has smaller
bad-sequence rank, so the induction hypothesis applies to it.

If $\fot(B)=0$, there are no moves from $B$, and \eqref{eq:residual} gives
$\fot(A+B)=\fot(A)$. Otherwise the contributions from moves in $B$ have
supremum
\[
 \sup_{x\text{ with a friend in }B}
 \bigl(\fot(A)+\fot(B\not\geq x)+1\bigr)
 =\fot(A)+\fot(B).
\]
Here associativity and continuity of ordinal addition in its right argument
justify the equality; when the supremum is a successor, it is attained.
Contributions from moves in $A$ have supremum $\fot(A)$ and cannot enlarge
this value. This proves \eqref{eq:binary-ordinal}.
\end{proof}

\begin{theorem}[Transfinite sum law]\label{thm:transfinite}
Let $\alpha$ be an ordinal and let $(P_\xi)_{\xi<\alpha}$ be wpos. Their
well-ordered sum is a wpo and
\begin{equation}\label{eq:transfinite}
 \boxed{\fot\left(\sum_{\xi<\alpha}P_\xi\right)
       =\sum_{\xi<\alpha}\fot(P_\xi).}
\end{equation}
\end{theorem}
\begin{proof}
In a bad sequence, the indices of its successive terms must be nonincreasing:
a term in a strictly earlier block is below every term in a later block.
An infinite nonincreasing sequence of ordinals is eventually constant, so
an infinite bad sequence in the sum would eventually give an infinite bad
sequence inside one block. Hence the sum is a wpo.

Prove the formula by transfinite induction on $\alpha$. The zero case is
immediate and successor stages are \cref{lem:binary-ordinal}. Let $\alpha$
be a nonzero limit. Every nonempty open-ended sequence starts in some block
$P_\xi$, and all its later choices lie in blocks of index at most $\xi$.
Every friend of a choice is in the same block as that choice. Thus the
entire subtree below this first choice is unchanged if the whole order is
replaced by its prefix through $P_\xi$.

Applying the rank recursion at the root therefore gives
\[
 \fot\left(\sum_{\xi<\alpha}P_\xi\right)
 =\sup_{\beta<\alpha}\fot\left(\sum_{\xi<\beta}P_\xi\right).
\]
The induction hypothesis and the definition of a limit ordinal sum turn
the right-hand side into $\sum_{\xi<\alpha}\fot(P_\xi)$.
\end{proof}

\begin{corollary}[Finite blocks]\label{cor:finite-blocks}
For any well-ordered family of finite posets,
\begin{equation}\label{eq:finite-blocks}
 \fot\left(\sum_{\xi<\alpha}P_\xi\right)
 =\sum_{\xi<\alpha}\bigl(|P_\xi|-\cc(\Inc(P_\xi))\bigr).
\end{equation}
In particular, if every incomparability component of a wpo $P$ is finite
and the components are $(C_\xi)_{\xi<\alpha}$ in their canonical order,
then
\begin{equation}\label{eq:finite-components-infinite}
 \fot(P)=\sum_{\xi<\alpha}(|C_\xi|-1).
\end{equation}
\end{corollary}
\begin{proof}
Use \cref{thm:main} in \cref{thm:transfinite}. For the second statement,
\cref{lem:ordered-components} provides the well-ordered component decomposition.
\end{proof}

The finite subtraction takes place \emph{inside each summand}; the outer sum
is ordinal addition in the specified order. Consequently a final finite
block can survive in the answer when the same block at the beginning would
be absorbed.

\begin{example}[Explicit ordinal widths]
Let $H=\sum_{n<\omega}\A_2$. Then
\[
 \fot(H)=\omega,\qquad w(\M_{\DM}(H))=\omega^\omega.
\]
For $k<\omega$,
\[
 \fot(H+\A_{k+1})=\omega+k,
 \qquad w(\M_{\DM}(H+\A_{k+1}))=\omega^{\omega+k}.
\]
Likewise, $\sum_{\xi<\omega^2}\A_2$ has friendly order type $\omega^2$
and multiset ordinal width $\omega^{\omega^2}$. Inserting arbitrary chain
blocks does not add to the friendly order type, since their contribution is zero.
\end{example}

\begin{corollary}[Every ordinal at ordinary width two]\label{cor:all-ordinals}
For every ordinal $\alpha$, put
\[
 P_\alpha=\sum_{\xi<\alpha}\A_2.
\]
Then $\fot(P_\alpha)=\alpha$ and
$w(\M_{\DM}(P_\alpha))=\omega^\alpha$. For $\alpha>0$, the ordinary
maximum antichain size of $P_\alpha$ is exactly two.
\end{corollary}
\begin{proof}
Every two-element antichain contributes one to \eqref{eq:transfinite}.
Antichains cannot meet two different blocks, and each nonempty block has
an antichain of size two. Apply the width-transfer theorem.
\end{proof}

\begin{proposition}[The infinite graph-only obstruction]\label{prop:graph-obstruction}
On infinite wpos, friendly order type is not determined by the isomorphism
type of the incomparability graph, even if every component is finite.
In fact, for all countably infinite ordinals $\alpha$, the graphs
$\Inc(P_\alpha)$ are isomorphic, while $\fot(P_\alpha)=\alpha$.
\end{proposition}
\begin{proof}
Each graph is a disjoint union of countably infinitely many isolated edges.
Any bijection between their block index sets induces a graph isomorphism.
The friendly order types are given by \cref{cor:all-ordinals} and need not
coincide; for example, $\omega$ and $\omega+1$ give different answers.
\end{proof}

Thus the finite graphic-rank interpretation is exact but does not extend
unchanged to infinite graphs. The component order supplies indispensable
ordinal information. The transfinite examples above are consequences of
the sum law, not claims that the general connected infinite case has been solved.

\section{Sharp finite bounds and loss under added comparisons}
\label{sec:bounds}
Let $\str(P)$ be the induced subposet obtained by deleting all vertices
that are comparable with every element of $P$. These are exactly the
isolated vertices of $\Inc(P)$.

\begin{proposition}[Exact stripped formula and sharp bounds]\label{prop:stripped}
Let $s=|\str(P)|$. Then $\fot(P)=\fot(\str(P))$. If $s=0$, the value is
zero. Otherwise $s\geq2$, and if $t$ is the number of non-singleton
incomparability components,
\[
 \fot(P)=s-t,\qquad
 \left\lceil\frac{s}{2}\right\rceil\leq\fot(P)\leq s-1.
\]
For each fixed $s\geq2$, every integer in this interval is attained.
\end{proposition}
\begin{proof}
Deleting an isolated vertex reduces the vertex count and the component
count by one, leaving their difference unchanged. The $t$ remaining
components have at least two vertices each, so $1\leq t\leq\lfloor s/2\rfloor$.
This proves the bounds. Conversely, for any such $t$, partition $s$ into
$t$ positive parts each at least two and take the ordinal sum of antichains
of those sizes. Its non-singleton component count is $t$.
\end{proof}

For an $n$-element poset with no restriction on isolated vertices, every
integer from $0$ through $n-1$ is possible: to realize $k$, use
$\C_{n-k-1}+\A_{k+1}$, with the empty-chain convention when necessary.

\begin{proposition}[Maximal separation at fixed elementary invariants]
For every $n\geq3$, the posets
\[
 L_n=\C_{n-2}+\A_2,
 \qquad R_n=\C_{n-1}\sqcup\C_1
\]
both have cardinality and maximal order type $n$, height $n-1$, and ordinary
width $2$, but
\[
 \fot(L_n)=1,\quad\fot(R_n)=n-1,
 \qquad
 w(\M_{\DM}(L_n))=\omega,\quad
 w(\M_{\DM}(R_n))=\omega^{n-1}.
\]
The friendly-order-type gap $n-2$ is the largest possible between two
non-chain $n$-element posets.
\end{proposition}
\begin{proof}
The chain summand of $L_n$ contributes $n-2$ isolated vertices and its
antichain contributes one component; hence $\fot(L_n)=1$. The singleton
in the disjoint union $R_n$ is incomparable to all $n-1$ chain elements,
so its incomparability graph is connected and $\fot(R_n)=n-1$.
The height and width assertions are immediate from the two constructions.
Every non-chain finite poset has friendly order type at least one by
\cref{lem:elementary} and at most $n-1$ by \eqref{eq:rough-bound}; hence
the gap is maximal.
\end{proof}

\begin{proposition}[Exact monotonicity under augmentation]
Suppose $P$ and $Q$ are finite posets on the same support and every
comparison of $P$ is also a comparison of $Q$. Then
\[
 \fot(P)-\fot(Q)=\cc(\Inc(Q))-\cc(\Inc(P))\geq0.
\]
\end{proposition}
\begin{proof}
Adding comparisons removes incomparability edges. Removing edges can only
split components, not join them. The equality is obtained by subtracting
the two instances of \cref{thm:main}.
\end{proof}
This statement applies to the completed transitive order. Adding one proposed
comparison can force many others by transitivity, so it need not lower the
invariant by only one.

\section{Enumeration and a bivariate generating function}
\label{sec:enumeration}
Call a poset on $\{1,\ldots,n\}$ \emph{naturally labelled} when $i<_P j$
implies $i<j$ as integers. Let $a_n$ count such posets, with $a_0=1$, and
let $b_n$ count the ones with connected incomparability graph for $n\geq1$.
These are not counts of unlabelled posets or of all labelled posets.
Every isomorphism type has a natural labelling, but may have many of them.

Put
\[
 A(z)=\sum_{n\geq0}a_nz^n,
 \qquad B(z)=\sum_{n\geq1}b_nz^n.
\]
These are formal power series; no analytic convergence is asserted.
Let $a_{n,k}$ count naturally labelled $n$-element posets with friendly
order type $k$.

\Needspace{16\baselineskip}
\begin{theorem}[Distribution by component rank]\label{thm:gf}
One has
\begin{equation}\label{eq:gf-univariate}
 A(z)=\frac{1}{1-B(z)},\qquad
 b_n=a_n-\sum_{i=1}^{n-1}b_i a_{n-i},
\end{equation}
and the bivariate distribution is
\begin{equation}\label{eq:gf}
 \boxed{\sum_{n,k\geq0}a_{n,k}z^n u^k
 =\frac{1}{1-\sum_{m\geq1}b_mz^m u^{m-1}}.}
\end{equation}
For $n\geq1$ and $0\leq k\leq n-1$,
\begin{equation}\label{eq:coefficient}
 a_{n,k}=[z^n]B(z)^{n-k}.
\end{equation}
\end{theorem}
\begin{proof}
In a natural labelling, if $C<D$ are incomparability components, every
label in $C$ is smaller than every label in $D$. Therefore components
occupy consecutive intervals of labels in their canonical order. The
component decomposition is consequently a unique ordered list of nonempty
naturally labelled connected-incomparability posets. Conversely, concatenating
any such list by ordinal sum gives a naturally labelled poset with exactly
those components. Summing over ordered lists gives
$1+B+B^2+\cdots$, proving the first identity. Extracting its coefficients
proves the recurrence.

A component of size $m$ contributes $m-1$ to friendly order type by
\cref{thm:main}. Weighting it by $z^m u^{m-1}$ gives \eqref{eq:gf}.
A poset of size $n$ and friendly order type $k$ has exactly $n-k$ components,
so extracting from lists of that length gives \eqref{eq:coefficient}.
\end{proof}

The exhaustive counts give
\[
 (a_0,\ldots,a_7)=(1,1,2,7,40,357,4824,96428)
\]
and
\[
 (b_1,\ldots,b_7)=(1,1,4,27,275,4070,86278).
\]
These values in this report are outputs of the included exact enumeration,
not numbers taken from a database without verification.

\begin{corollary}[Fixed-rank counting polynomials]\label{cor:polynomial}
For fixed $k\geq0$ and $n\geq k+1$,
\begin{equation}\label{eq:fixed-rank}
 a_{n,k}=
 \sum_{\substack{m_1,\ldots,m_k\geq0\\\sum_{j=1}^k j m_j=k}}
 \frac{(n-k)_{\underline r}}{\prod_{j=1}^k m_j!}
 \prod_{j=1}^k b_{j+1}^{\,m_j},
 \qquad r=\sum_{j=1}^k m_j,
\end{equation}
where $(x)_{\underline r}=x(x-1)\cdots(x-r+1)$ and empty products are $1$.
It is a polynomial in $n$ of degree $k$ with leading coefficient $1/k!$.
\end{corollary}
\begin{proof}
There are $n-k$ component positions. Let $m_j$ record how many components
have size $j+1$. Their total rank contribution is $\sum jm_j=k$, while
all remaining components are singletons. Choose and order the $r$ non-singleton
positions in $(n-k)_{\underline r}/\prod m_j!$ ways, then choose the connected
structures on each position. This gives the formula.
The degree is at most $r\leq k$. Degree $k$ occurs only when $m_1=k$;
since $b_2=1$, that term is $\binom{n-k}{k}$ and has leading coefficient $1/k!$.
\end{proof}

For example, in their indicated range of $n$,
\begin{align*}
 a_{n,0}&=1,\\
 a_{n,1}&=n-1,\\
 a_{n,2}&=4(n-2)+\binom{n-2}{2}
          =\frac{(n-2)(n+5)}2,\\
 a_{n,3}&=27(n-3)+4(n-3)(n-4)+\binom{n-3}{3}.
\end{align*}
These formulas offer a separate combinatorial explanation for the initial
columns of the computation table, rather than only an empirical pattern.

\section{Algorithms and certificate verification}
\label{sec:algorithms}
The package implements two intentionally different ways to calculate the
finite invariant. The structural algorithm constructs the incomparability
graph, counts connected components, and returns $n-c$. The reference
algorithm uses only the defining residual recursion.

\subsection{Direct computation of the value}
With a complete $n\times n$ comparison matrix, incomparability testing and
graph construction take $O(n^2)$ Boolean operations. Connected components
can be found in $O(n^2)$ time in this dense representation. Thus the friendly
order type and the finite exponent in the multiset-width formula are
computed in $O(n^2)$ time from a valid comparison matrix. The returned
ordinal is represented symbolically as $\omega^k$; there is no attempt to
enumerate the infinite multiset order.

Input validation is separate. A naive transitivity check takes $O(n^3)$
Boolean operations, and input consisting only of generating strict relations
requires transitive closure first. The supplied command-line interface
accepts such relations, computes their closure, and rejects cycles. The
quadratic claim is for evaluation after a valid comparison matrix is given,
not for all possible input representations.

\subsection{Constructing an optimal witness}
Decompose the graph into its ordered components. Process them in decreasing
component order. Within a component, repeatedly apply \cref{lem:noncut},
record an incomparable neighbor as friend, and delete the selected maximal
vertex. Stop when only the chosen root remains. A root can be specified
in advance, provided it is minimal in that component.

An implementation of one non-cut step first selects a maximal vertex,
computes the components after its deletion, and either uses that vertex or
selects a maximal element of the last new component as in the proof.
Each such step needs $O(n^2)$ elementary operations on a comparison matrix,
so the full witness can be constructed in $O(n^3)$ time. It contains at most
$n-1$ choice--friend pairs. The Python implementation uses integer bit sets;
these matrix bounds describe the elementary-operation model, rather than
pretending that arbitrarily long machine integers have unit cost.

A verifier simulates the residuals and checks that each named choice and
friend is still present and incomparable. It also checks that the witness
edges are acyclic and span each original component. This establishes
legality and checks the claimed forest certificate. Optimality is then
certified by comparing the sequence length with the component upper bound.

\subsection{Independent residual recursion}
For a subset mask $R$, define
\begin{equation}\label{eq:dp}
 F(R)=\max\left(\{0\}\cup
 \{1+F(R\setminus\mathord\uparrow x):x\in R,\ R\cap(P\perp x)\ne\varnothing\}\right).
\end{equation}
Memoizing this recurrence considers at most $2^n$ subsets, with at most $n$
choices at each subset. Its worst-case count is $O(n2^n)$ subset transitions,
before accounting for the cost of operations on the bit sets. It does not
call the component formula and does not use the maximal non-cut lemma.
Thus comparison against it is not merely testing the formula against itself.

\subsection{Using the supplied program}
The example input \texttt{examples/diamond.json} is
\begin{lstlisting}
{"n":4,"relations":[[0,1],[0,2],[1,3],[2,3]]}
\end{lstlisting}
Run the following from the archive directory:
\begin{lstlisting}
python src/friendly.py examples/diamond.json --dp
python src/verify.py --max-n 7
\end{lstlisting}
The first command reports components $\{0\},\{1,2\},\{3\}$, friendly order
type $1$, multiset ordinal width \texttt{omega\textasciicircum(1)}, and an independently
checked certificate. Input labels are integers from $0$ through $n-1$;
they need not themselves form a linear extension. Optional input
\texttt{"roots":[...]} prescribes one minimal root per component. The code
uses Python 3.10 or later and only the standard library.

\section{Exact computational checks}
\label{sec:verification}
The delivered run was performed with Python 3.13.5. It generated every
naturally labelled poset through size seven. To pass from size $n-1$ to
size $n$, the strict predecessor set of the new largest label is chosen
among the order ideals of the previous poset. Every such ideal gives a
valid natural order, and every natural order arises once in this way.
Since every finite poset admits a linear extension, all isomorphism types
through the bound are represented, generally with multiplicity.

For each generated order, the test compares \eqref{eq:dp} with the graph
formula, repeats the direct recursion for the dual order, constructs an
optimal certificate, and simulates it. The maximal non-cut conclusion is
additionally checked at the root of every connected case of size at least two.

\begin{center}
\small
\begin{tabular}{r r r r r r r r r}
\toprule
$n$ & Total & $f=0$ & $f=1$ & $f=2$ & $f=3$ & $f=4$ & $f=5$ & $f=6$\\
\midrule
0 & 1 & 1 & -- & -- & -- & -- & -- & --\\
1 & 1 & 1 & -- & -- & -- & -- & -- & --\\
2 & 2 & 1 & 1 & -- & -- & -- & -- & --\\
3 & 7 & 1 & 2 & 4 & -- & -- & -- & --\\
4 & 40 & 1 & 3 & 9 & 27 & -- & -- & --\\
5 & 357 & 1 & 4 & 15 & 62 & 275 & -- & --\\
6 & 4,824 & 1 & 5 & 22 & 106 & 620 & 4,070 & --\\
7 & 96,428 & 1 & 6 & 30 & 160 & 1,047 & 8,906 & 86,278\\
\bottomrule
\end{tabular}
\end{center}
The total is 101,660 orders, with 90,655 explicit connected-case non-cut
checks. Additional tests cover all 2,601 ordered pairs of naturally labelled
factors of sizes zero through four, comparing both graph counts and the
independent residual recursion with the Cartesian-product formula. The
largest product supports have sixteen elements.

For substitution, every natural base of size at most four is tested with
each fiber independently selected from a singleton, a two-element chain,
a two-element antichain, and a three-element chain. This gives 10,725
heterogeneous substitutions, each checked against direct residual recursion.
For maximum friendly subsets, 2,365 candidate subsets of the predicted
maximum size are tested through support size five by a recursion restricting
choices to that subset. Exactly the subsets in \cref{thm:subsets} pass;
1,052 corresponding prescribed-root certificates are constructed and verified.
Eight malformed-order or malformed-friend cases are rejected as expected.

All tests passed. The full run recorded 15.584 seconds in this environment;
this is a reproducibility datum, not a performance guarantee. The JSON
report and CSV rank distribution are included under \texttt{data/}.

The generator is exhaustive only in the stated finite ranges. It does not
verify a theorem for all finite posets, and it does not verify a transfinite
rank computation. The proofs above supply those assertions. The computation
is useful for detecting definition mistakes, counterexamples, off-by-one
errors, incorrect empty cases, and errors in certificate construction.
No Lean formalization is claimed.

\section{What has been settled, and what has not}
\label{sec:scope}
The finite specialization has an exact answer with a complete proof:
\[
 \fot(P)=|P|-\cc(\Inc(P)).
\]
It yields all maximum friendly subsets, efficient certificates, finite
substitution and product laws, sharp bounds, and an enumerative distribution.
Together with a sum law, it also evaluates all wpos whose incomparability
components are finite.

The broader problem from the literature remains larger than these statements.
A wpo may have infinite incomparability components. The connected finite
proof relies essentially on maximal vertices and finite deletion. Neither
is available in general. For example, an infinite wpo may have no maximal
vertex at all, so the finite construction is not a transfinite algorithm
in disguise. Moreover, \cref{prop:graph-obstruction} proves that even a
complete description of the abstract infinite incomparability graph cannot
suffice unless its component order is retained.

A concrete remaining direction is therefore to understand the friendly ranks
of \emph{infinite connected} wpos and how they respond to finite extensions.
The component decomposition reduces a general instance to those connected
pieces plus ordinal addition, but it does not calculate their ranks. We
make no conjectural formula for those pieces on the basis of the finite
experiment alone.

The proposed new finite characterization warrants external checking and a
broader literature comparison. If it is already implicit in an earlier
source, the proofs and artifacts remain useful but the novelty claim would
need to be revised. This distinction is separate from the question whether
the displayed theorems follow from the arguments given.

\appendix
\Needspace{23\baselineskip}
\section{Statement and dependency audit}
\label{sec:audit}
The following table separates imported results from the mathematical work
proved in this manuscript.
\begin{center}
\small
\begin{tabular}{p{0.34\textwidth} p{0.57\textwidth}}
\toprule
Statement & Dependency and status\\
\midrule
Definition of friendly order type & Vialard's invariant; not introduced here.\\
Multiset-width transfer & Imported theorem \eqref{eq:transfer}; no new proof claimed.\\
Ordered components & Elementary order argument, proved in \cref{lem:ordered-components}.\\
Protected-root non-cut lemma & Full finite proof in \cref{lem:noncut}.\\
Finite component formula & Capacity bound plus constructive lower bound; \cref{thm:main}.\\
Maximum friendly subsets & Protected-root construction and necessity argument; \cref{thm:subsets}.\\
Finite multiset-width formula & Component theorem plus imported width transfer.\\
Binary ordinal-sum law & Known law, independently proved in \cref{lem:binary-ordinal}.\\
Transfinite finite-block formulas & Ordinal-sum rank proof plus finite theorem.\\
Product and substitution laws & Direct graph arguments plus finite theorem.\\
Counting generating function & Unique ordered component decomposition of natural labellings.\\
Finite computational agreement & Executed tests within explicit bounds; not a formal proof.\\
Novelty/priority & Not certified by literature search or by the computation.\\
\bottomrule
\end{tabular}
\end{center}

\subsection*{Checks that prevent common misstatements}
The friend must be in the current residual, isolated vertices must be counted
as components, and an empty base poset gives a singleton multiset order of
width one. The product formula requires both factors to have at least two
elements. Empty substitution fibers require deleting their base vertices
first. Protected roots must be minimal \emph{inside their component}, not
necessarily globally minimal in the original poset. Forest roots for a later
component may be removed when an earlier component is processed; they remain
roots of the certificate forest, not necessarily members of the final residual.
Finally, the outer addition in a transfinite sum is ordinal addition, not
natural addition or cardinal addition.

\section{Archive contents and reproduction}
The main files are \texttt{article.tex}, \texttt{article.pdf},
\texttt{sources.bib}, and \texttt{README.md}. Executable code is in
\texttt{src/friendly.py} and \texttt{src/verify.py}; concrete inputs are in
\texttt{examples/}; verification output is in \texttt{data/}. The accompanying
\texttt{build.sh} compiles the article using a standard TeX installation.
No font files, downloaded source papers, checksum files, or claims of external
peer review are included.

The source and proofs are intentionally separable. The component formula
can be checked by reading only \cref{sec:definitions,sec:components,sec:main};
its finite applications do not depend on transfinite notation. The ordinal
consequences require \cref{sec:transfinite}, and the width consequences also
require the explicitly imported theorem. The independent dynamic program
is confined to finite computations and is not used as a premise in any proof.

\Needspace{24\baselineskip}
\bibliographystyle{plain}
\bibliography{sources}
\end{document}
